% Theme: Chalkboard  --  best for derivations, built up one line per swipe
% Build:  ./build.sh example-chalkboard.tex   (or compile on Overleaf)
\documentclass[12pt]{beamer}
\usetheme{igchalkboard}

\ighandle{@quantum_danesh}
\igseries{Qubits desde cero}
\igissue{05}

\begin{document}

\coverslide{Pizarra}
  {Despolarizar un qubit, línea a línea}
  {Un canal, tres líneas de álgebra y un dibujo.}

\begin{frame}{El planteamiento}{Datos}
Todo estado de un qubit es un punto de la bola de Bloch:
\[
  \rho = \frac12\big(\id + \Vec r\cdot\Vec\sigma\big),
  \qquad |\Vec r| \le 1 .
\]
El canal despolarizante conserva $\rho$ con probabilidad $1-p$ y lo
cambia por ruido puro con probabilidad $p$:
\[
  \mathcal{D}_p(\rho) = (1-p)\,\rho + p\,\frac{\id}{2}.
\]
\key{Pregunta:} ¿qué le hace al vector de Bloch $\Vec r$?
\end{frame}

\begin{frame}{Sustituimos}{Derivación}
{\large
\begin{align*}
  \mathcal{D}_p(\rho)
  &= (1-p)\,\frac12\big(\id + \Vec r\cdot\Vec\sigma\big) + \frac{p}{2}\,\id\\
  &\uncover<2->{= \frac12\big[(1-p) + p\big]\,\id + \frac12(1-p)\,\Vec r\cdot\Vec\sigma}\\
  &\uncover<3->{= \frac12\big(\id + \mhl<4->{(1-p)\,\Vec r}\cdot\Vec\sigma\big)}
\end{align*}}%
\uncover<4->{Solo cambia el vector de Bloch:
\keyeq{\Vec r \;\longmapsto\; (1-p)\,\Vec r}}
\end{frame}

\begin{frame}{Ahora, iteramos}{Punto fijo}
\[
  \Vec r_n = (1-p)^n\,\Vec r_0 \;\xrightarrow[\;n\to\infty\;]{}\; \Vec 0 .
\]
Todo estado fluye hacia el centro de la bola, $\id/2$. Como $1-p<1$, el
mapa es una \hl{contracción} y el teorema de Banach garantiza que este
punto fijo es único.
\begin{center}
\begin{tikzpicture}[x=1mm,y=1mm]
  \pgfmathsetseed{8}
  \draw[chalk, igFG, line width=0.7pt] (0,0) circle (17.5);
  \draw[chalk, igFG, opacity=0.5, line width=0.4pt] (0,0) ellipse (17.5 and 4.6);
  \foreach \n/\s in {0/1, 1/0.6, 2/0.36, 3/0.216}{
    \fill[igAccent] ({17.5*\s*cos(35)},{17.5*\s*sin(35)}) circle (0.8);
  }
  \draw[chalk, igAccent, line width=0.5pt, -{Latex[length=1.6mm]}]
    ({17.5*cos(35)},{17.5*sin(35)}) -- ({17.5*0.24*cos(35)},{17.5*0.24*sin(35)});
  \node[igAccent, anchor=west] at ({17.5*cos(35)+1},{17.5*sin(35)+1}) {$\Vec r_0$};
  \node[igAccent, anchor=north west] at ({10.5*cos(35)+1},{10.5*sin(35)-0.5}) {$\Vec r_1$};
  \fill[igAccentB] (0,0) circle (0.9);
  \node[igAccentB, anchor=north] at (0,-1.3) {$\id/2$};
\end{tikzpicture}
\end{center}
\end{frame}

\begin{frame}{Comprobamos los Kraus}{Verificación}
\[
  K_0 = \sqrt{1-\frac{3p}{4}}\;\id ,
  \qquad
  K_j = \frac{\sqrt{p}}{2}\,\sigma_j
\]
\[
  \sum_k K_k^\dagger K_k
  = \Big(1-\frac{3p}{4}\Big)\id + 3\cdot\frac{p}{4}\,\id
  = \mhl{\id}
\]
para $j = x,y,z$. Preserva la traza. \aside{\checkmark}
\begin{notebox}
Con $\sum_j \sigma_j\rho\,\sigma_j = 2\,\id - \rho$ para $\Tr\rho = 1$,
la suma de Kraus devuelve $(1-p)\rho + p\,\id/2$.
\end{notebox}
\end{frame}

\begin{frame}{Te toca}{Tarea}
El amortiguamiento de amplitud tiene
$K_0 = \mathrm{diag}\big(1,\sqrt{1-\gamma}\big)$ y
$K_1 = \sqrt{\gamma}\,\ket{0}\!\bra{1}$. Demuestra que mueve el vector de
Bloch así:
\[
  \begin{pmatrix} r_x \\ r_y \\ r_z \end{pmatrix}
  \longmapsto
  \begin{pmatrix} \sqrt{1-\gamma}\, r_x \\ \sqrt{1-\gamma}\, r_y \\
    \mhl*{\gamma} + (1-\gamma)\, r_z \end{pmatrix}
\]
\hl{¿Dónde termina cualquier estado?}
\end{frame}

\closingslide{Hasta aquí la pizarra}
  {La respuesta de la tarea, en los comentarios. El post completo está en el blog.}

\end{document}
